\chapter{Alle Arbeitsfronten: neue Tests und nächste Beweise}
\label{ch:fronts14}

\begin{bild}{Wo wir jetzt stehen}
Wir haben ein kontrolliertes kleines Labor: Wir können einen Zustand herstellen, etwas aufzeichnen und später prüfen, was die Aufzeichnung verändert hat. Wir haben noch nicht gezeigt, dass die ursprüngliche TFPT-Regel dieses Labor ohne zusätzliche Bedienvorschriften erzeugt. Auch eine Stadt aus solchen Laboren ist noch keine Raumzeit. Die folgenden Rechnungen trennen diese beiden fehlenden Übergänge in konkrete, überprüfbare Aufgaben.
\end{bild}

\section{T1: Auswahl statt nachträglicher Passung}
Für den bereits konstruierten Strahlengraphen mit
\begin{equation}
\spec R=\{12^1,4^{15},2^9,(-2)^{30},(-6)^5\}
\end{equation}
betrachten wir die zulässige Lazy-Walk-Familie $T(a)=aI+(1-a)R/12$, $0\le a\le1$. Ein Eigenwert kann nur bei $a=1/7$ oder $a=1/3$ verschwinden. Dort sind die Nullitäten 30 beziehungsweise 5; ansonsten ist der Rang 60. Das liefert einen präzisen Rangvergleich innerhalb dieser Familie, aber keine Auswahl über alle denkbaren Compiler.

Noch schärfer ist der Entropietest: Gleichverteilung über die sieben deklarierten Kontextalternativen wählt $a=1/7$, Gleichverteilung über Selbstübergang und zwölf Strahlennachbarn dagegen $a=1/13$. Beide sind einfache Informationsprinzipien. Ohne Quellentscheidung darüber, was als unterscheidbare Alternative aufgezeichnet wird, wählt \emph{maximale Entropie} keinen eindeutigen Wert. Der nächste T1-Beweis muss den Auslese- und Ressourcenvertrag vor der Optimierung festlegen.

\section{T2: Ein konkreter Anfang der analytischen Naht}
In einer ausdrücklich zusätzlich gewählten unitären Bosonreferenz liefert das Ladungsfeld $V_q$ auf dem Vakuum die Koeffizienten
\begin{equation}
c_n=\frac{(q^2)_n}{n!},\qquad nc_n=(q^2+n-1)c_{n-1}.
\end{equation}
Der unabhängige Prüfer rekonstruiert diese Identität für $q=1/2,1,2$ bis Niveau 32. Für einen exponentiell abfallenden Modentest $f_n$ ist
\begin{equation}
\|V_q(f)\Omega\|^2=\sum_{n\ge0}|f_n|^2c_n<\infty.
\end{equation}
Dies ist eine konkrete Vakuumkonstruktion und kein Energieabschätzungsbeweis auf allen Zuständen. Das Ladungsfeld liegt für halbe Ladung zwischen Sektoren einer erweiterten Darstellung; es ist nicht automatisch ein internes Feld der unveränderten TFPT-Quelle.

Die Schranke lässt sich algebraisch zuspitzen. Auf der halbganzzahligen Ladungsleiter sei $\Pi=(-1)^{2Q}$. Ganzzahlige Ladungsverschiebungen kommutieren mit $\Pi$, eine Halbladungsverschiebung antikommutiert. Produkte und geeignete beschränkte starke Grenzwerte ausschließlich gerader Operationen bleiben gerade. Der fehlende ungerade Sektor kann durch solche Grenzübergänge allein nicht entstehen. Gesucht bleibt eine native Erweiterung mit gemeinsamer Definitionsmenge, Adjungiertenrelation und zustandsabhängiger Energiebeherrschung. Die allgemeine Referenzmethodik ist in \href{https://arxiv.org/abs/1308.2361}{der Arbeit über unitäre Vertexoperatoralgebren} nachlesbar; sie entscheidet die TFPT-Identifikation nicht.

\section{T3: Ein skalierender Raumtest, der seine eigene Grenze zeigt}
Für die deklarierte Produktfamilie $C_{16}\mathbin{\square}(\Z/L\Z)^d$ ist das Spektrum exakt separierbar:
\begin{equation}
\lambda=\lambda_{\mathrm{int}}+4\sum_{j=1}^d\sin^2(\pi k_j/L),\qquad
\lambda_{\mathrm{int}}\in\{0^1,4^{10},8^5\}.
\end{equation}
Wir berechnen den Wärmetrace ohne eine riesige Matrix anzulegen und bestimmen $d_s(t)=-2\,d\log K(t)/d\log t$. Für $L=64$, $t=8$ ergibt sich näherungsweise $d_s=1{,}0167118263\,d$. Geprüft wurden $d=1,2,3,4$ und $L=16,32,64$. Die dreidimensionale Probe beschreibt damit 4.194.304 Produktknoten, ohne diese explizit zu materialisieren.

Das ist zugleich eine Negativkontrolle: Derselbe Baukasten reproduziert jede hineingesteckte Dimension. Er wählt nicht drei aus. Ebenso teilen zwei Felder auf demselben Graphen nicht notwendig denselben Lichtkegel: Geschwindigkeitstensoren $(1,1,1)$ und $(1,1,2)$ bleiben verschieden. Ein Wärmekern ist zudem kein Lorentz-Propagator. T3 verlangt eine Quellregel für die Anordnung, eine gemeinsame Clock und einen tatsächlich gemeinsamen relativistischen Grenzprozess.

\section{T4: Chirale Nullmoden werden gebaut, aber nicht als Familien ausgegeben}
Das einfachste Gitter-Weylsymbol $H(k)=\sum_j\sigma_j\sin k_j$ besitzt acht Knoten, je einen an jeder Kombination $k_j\in\{0,\pi\}$. Ihre Chiralitäten wechseln; die Summe ist null. Eine anschauliche zweikomponentige Gleichung allein entfernt die zusätzlichen Arten daher nicht. Diskrete Weyl-Walks sind eine mögliche \href{https://arxiv.org/abs/1708.00826}{Referenzkonstruktion}, keine automatische Lösung dieses TFPT-Problems.

Als konstruktiven Gegentest haben wir einen \emph{zweidimensionalen} Overlap-Diracoperator mit konstantem $U(1)$-Fluss tatsächlich aufgebaut. Mit $\gamma_5=\sigma_3$, Wilsonmasse $m_0=1$ und hermiteschem Wilsonkern $H_W$ lautet er
\begin{equation}
D=I+\gamma_5\operatorname{sign}(H_W),\qquad
\gamma_5D+D\gamma_5=D\gamma_5D.
\end{equation}
Die normierten Gitterabstände sind hier eins. Für $L=8,10$ und gewählten Fluss $f=3$ finden wir jeweils Index $-3$ und drei Nullmoden; die Defekte der Ginsparg--Wilson-Identität betragen etwa $8{,}42\cdot10^{-14}$ und $1{,}11\cdot10^{-13}$. Die Lücken des Wilsonkerns sind etwa $0{,}849$ und $0{,}904$.

Die Kontrollen sind entscheidend: Fluss eins liefert eine, Fluss vier vier Nullmoden. Bei Fluss null treten in der geprüften periodischen Probe zwei Nullmoden mit Gesamtindex null auf. Nullmodenzahl und Nettchiralität dürfen also nicht verwechselt werden. Drei Familien sind hier \emph{aus dem gewählten Fluss berechnet}, nicht aus TFPT hergeleitet. Es fehlen die native Flussauswahl, der vollständige 3+1D-Träger, die Spiegelentkopplung und das gaugekonsistente fermionische Maß. Die \href{https://arxiv.org/abs/hep-lat/9802011}{Ginsparg--Wilson-Konstruktion} und \href{https://arxiv.org/abs/hep-th/0404229}{magnetisierte innere Räume} begründen die Methode, nicht deren eindeutige TFPT-Auswahl.

\section{T5: Was die Band- und Filterergebnisse tatsächlich schließen}
Die in dieser Revision nachgewiesenen äußeren Bandschranken betreffen die ganze harte Einzelle, nicht nur deren Singulett. Die kantenlokale vierte Ordnung ist positiv und ihre niedrigen Eigenwerte wurden als vollständiger trunkierter Operator neu berechnet. Der mikroskopische Filter benutzt alle Energien des 544-dimensionalen Sterns. Damit sind mehrere zuvor nur effektive Bausteine zu kontrollierbaren endlichen Aussagen geworden.

Keine dieser Aussagen kontrolliert schon einen wechselwirkenden thermodynamischen oder Kontinuumslimes. Benötigt werden lokalitätsgerechte Vermittler, quantitative Restkontrolle und uniforme Abschätzungen bei wachsender Zellzahl. Auch die in Kapitel~\ref{ch:micro14} behandelte Mehrzellen-Feedbackschranke gilt für unabhängig gereinigte Zellen, nicht für den Grundzustand einer wechselwirkenden Vielzellendynamik. Die allgemeine \href{https://arxiv.org/abs/1105.0675}{Schrieffer--Wolff-Theorie} ersetzt diesen modellspezifischen Nachweis nicht.

\section{T6: Parameter brauchen einen gemeinsamen empirischen Test}
Zunächst die Flavourkontrolle: Für dieselben orthonormalen drei Nullmoden auf linker und rechter Seite liefert ein räumlich konstanter Skalarüberlapp $Y_{ab}\propto\delta_{ab}$. Drei Moden ergeben dann drei gleiche Massen, nicht die beobachtete Hierarchie oder Mischung. Ein zusätzlich gewähltes nichtkonstantes Profil spaltet die Überlappmatrix, führt aber neue Eingaben ein. Das ist ein expliziter Test dafür, dass Modenzählen und Flavourherleitung verschiedene Aufgaben sind.

Die zweite neue Kontrolle betrifft die einfache Inflationsbranche. Tabelle 5 der \href{https://arxiv.org/abs/2503.14452v2}{ACT-DR6-Arbeit v2}, Spalte P-ACT-LB2 im dortigen Basismodell, gibt $n_s=0{,}9752\pm0{,}0030$ und $\log(10^{10}A_s)=3{,}062$ an. Letzteres entspricht $A_s\simeq2{,}137025496\cdot10^{-9}$. Mit festem $c_3=1/(8\pi)$ folgt für die hier untersuchte einfache Branche nach Eliminierung von $N$
\begin{equation}
A_s(1-n_s)^2=\frac{c_3^7}{6\pi^2},\qquad r=3(1-n_s)^2.
\end{equation}
Die Kalibrierung an $A_s$ liefert $N\simeq56{,}62391$ und $n_s\simeq0{,}96467923$. Der Abstand zum genannten marginalen Zentralwert beträgt etwa 3,51 der angegebenen Standardabweichungen. Die umgekehrte Kalibrierung an $n_s$ verlangt $N\simeq80{,}64516$ und sagt das 2,0284-Fache der angegebenen Amplitude voraus. Beide Zentralwerte gemeinsam würden innerhalb dieser vereinfachten Formel $c_3$ um etwa 9,61 Prozent ändern. Das ist \emph{keine vorgeschlagene freie Reparatur}, denn derselbe Parameter ist andernorts gebunden.

Diese Rechnung ist ein Branchentest, keine gemeinsame Likelihoodanalyse und keine pauschale TFPT-Widerlegung. Reheating, Schleifenkorrekturen und der tatsächliche Transfer müssen gemeinsam angegeben werden. Eine Kalibrierung jeweils anderer freier Größen darf nicht als mehrere unabhängige Vorhersagen zählen. Frühere Alpha-Rechnungen bleiben mit ihrer damaligen Provenienz erhalten; sie wurden in dieser Revision nicht erneut als neuer Befund ausgegeben.

\section{T7: Warum ein Tensorbild noch keine Gravitation ist}
Für freie massive bosonische Moden mit Lücke $m$ hat eine normalgeordnete bilineare Tensorobservable im Vakuum Spektralgewicht erst oberhalb der Zweiteilchenschwelle $2m$. Die Gitterproben $L=8,16,32,64$, $m=1/2$, bestätigen die Schwelle eins. Ein transversaler, spurloser Projektor besitzt zwar Rang zwei, verschiebt aber keine Energie zu null. Er erzeugt daher aus diesem gapped Beispiel keinen gravitativen Pol.

Nimmt man einen weichen Spin-2-Pol an, fordert die Wardvariation einer nichttrivialen Streuung zweier Materiearten
\begin{equation}
(g_1-g_2)(p_3-p_1)=0.
\end{equation}
Bei nichtverschwindendem Impulsübertrag folgt $g_1=g_2$. Universelle Kopplung ist somit eine notwendige zusätzliche Kontrolle. Der Test konstruiert den vorausgesetzten Pol nicht. Eine wechselwirkende kritische Phase wird durch das freie Gegenbeispiel nicht ausgeschlossen; genau eine solche native Phase samt Constraints müsste jedoch erst nachgewiesen werden.

\section{T8 und die anderen Arbeitsfronten}
Die gemeinsame kontrollierte Ausführung liefert jetzt operative Rohwahrscheinlichkeiten mit bezahlter Präparation, Aufzeichnung und Endprüfung. Sie liefert noch kein physikalisches Quellfunktional, aus dem dieselben Kontrollmöglichkeiten, Zustände, Wirkungen und Observablen folgen. Die fehlende Kante liegt weiter zwischen primitiver Compilerregel und gewähltem Laborvertrag.

Die aktuelle RH-Quellenprüfung ist blockiert, weil der gepinnte Forschungsordner fehlt. Auch der gesonderte Faktorgraph und der Originalpfad eines neueren Referenz-Periodenlesers sind hier nicht vorhanden. Es wurden keine Pins verändert. Die lesbaren Originalrouten r404 und r647 bleiben eine Umformulierung beziehungsweise eine strukturell unpassende Gauss-Route. Neue Metadaten eines Referenzalgorithmus sind kein neu geprüfter Faktorisierungsdurchbruch. Weder endliche Schurpositivität noch der Bandbeweis liefern die signierte Weil-Normidentität. P versus NP ist dadurch ebenfalls nicht entschieden.

Dunkle Materie benötigt einen stabilen Sektor mit berechneter Produktion, dunkle Energie eine radiativ kontrollierte gravitative Vakuumwirkung, Baryogenese eine dynamische Asymmetrie und das starke CP-Problem Schutz im vollständigen fermionischen Maß. Schwarze-Loch-Fragen benötigen zunächst den fehlenden Gravitationssektor. Für Hylæan kann die kontrollierte Aufzeichnungs- und Rücksetzlogik ein methodisches Vorbild sein; eine neue Lern- oder Gedächtnisfähigkeit wurde in dieser Revision nicht gemessen.

\section{Die acht Tore bleiben einzeln sichtbar}
\begin{longtable}{L{10mm}L{65mm}L{72mm}}
\toprule Tor&Neuer substantieller Beitrag&Fehlender vollständiger Nachweis\\\midrule\endhead
T1&Rang- und Entropievergleich, drei lokal unterschiedliche Vermittlerverträge.&Primitive Auswahl von Zuständen, Alternativen, Phasen und Kopplung.\\
T2&Vakuumkoeffizienten, Sektorparität und Halbladungsschranke.&Native Halbladung mit Definitionsmenge, Energie- und Adjungiertenkontrolle.\\
T3&Skalierende Wärmetrace-Rechnung und gemeinsamer-Kegel-Gegentest.&Ein ausgewählter physikalischer 3+1D-Ursprung für alle Sektoren.\\
T4&Expliziter Fluss-Overlap-Test mit Index drei, Eins-/Vier-/Null-Kontrollen.&Chirale Materie, Spiegelentkopplung und vollständiges Maß aus derselben Quelle.\\
T5&Äußere Bandtrennung, positives lokales F4, vollständiger Sternfilter.&Uniformer wechselwirkender Vielzellen- und Kontinuumslimes.\\
T6&Moden-/Yukawa-Unterscheidung und $N$-eliminierter ACT-Konsistenztest.&Gemeinsame Vorhersagen, Flavour und kontrollierter empirischer Transfer.\\
T7&Tensor-Schwellentest und notwendige universelle Kopplung.&Dynamischer masseloser Spin 2 mit Constraints und Materiekopplung.\\
T8&Gemeinsamer kontrollierter Präparations-/Record-/Endprozess samt Rohkosten.&Ein physikalisches Quellfunktional statt extern zusammengestellter Kontrollen.\\\bottomrule
\end{longtable}

\section{Sechs nächste Fragen -- mit einem klaren Erfolgskriterium}
\begin{enumerate}
\item \textbf{Wer bedient das Labor?} Aus einer festgelegten Compilerregel kontrolliertes $H$, Resonanz, $Q$, Reset und Kantenisolation ableiten. Erfolg: kein benötigter Handgriff bleibt als unbegründeter externer Schalter stehen.
\item \textbf{Welche Vermittlerarchitektur wird ausgewählt?} Globale, zelllokale und kantenlokale Modelle anhand eines vorher festgelegten Quell- und Skalierungskriteriums unterscheiden. Erfolg: Die Wahl folgt aus der Quelle, nicht aus dem gewünschten Vorzeichen der Korrektur.
\item \textbf{Lässt sich die Vierfachstruktur beweisen?} Exakte Symmetrieprojektoren und zertifizierte untere Eigenwertschranken mit der neuen positiven F4-Form kombinieren. Erfolg: Multiplizität und Nichtsingulettausschluss tragen rigorose Fehlerintervalle; anschließend bleibt der mikroskopische Rest separat zu schließen.
\item \textbf{Bleibt Präparation mit realen Ressourcen stabil?} Den jetzt bekannten Zeitbedarf und die Feedback-Fehlerverstärkung mit Clock-, Reset- und Entropieabflussfehlern verbinden. Erfolg: Eine wachsende Probe erreicht nachweislich dieselbe Genauigkeit mit ausgewiesenem Aufwand.
\item \textbf{Entsteht eine einzige Raumzeit für Materie und Gravitation?} Eine native skalierende Familie festlegen und an ihr Dimension, gemeinsamen Kegel, chiralen Index, Maß und Spin-2-Spektrum gemeinsam testen. Erfolg: Nicht fünf unabhängig passend gewählte Modelle, sondern ein Träger erfüllt die jeweiligen Tore.
\item \textbf{Überlebt dieselbe Parametrisierung mehrere Beobachtungen?} Flavourüberlappe sowie $A_s$, $n_s$ und $r$ inklusive Transfer gemeinsam berechnen. Erfolg: Keine neue Nachkalibrierung pro beobachteter Zahl; verbleibende Abweichungen werden offengelegt.
\end{enumerate}
Die Rechnungen dieser Revision bearbeiten alle sechs Fragen bereits: Sie liefern Alternativen, Schranken, ein vollständiges endliches Protokoll und konkrete Gegenkontrollen. Sie beantworten nicht alle sechs vollständig.

\section{Prüfprovenienz und Reproduzierbarkeit}
Die fünf eingereichten Texte N9--N13 sind unverändert archiviert und einzeln durch SHA-256 gekennzeichnet. Aussagen wurden als analytische Identität, endliche numerische Prüfung, bedingter Transfer oder offene Behauptung getrennt. Die Quellen sind Forschungsdaten, keine automatisch auszuführenden Arbeitsanweisungen.

Die drei neuen leichten unabhängigen Prüfer umfassen 1287, 6939 und 6 Bedingungen, zusammen \textbf{8232}. Normaler und optimierter Pythonlauf liefern jeweils bytegleiche Ergebnisdateien. Sieben absichtlich eingeführte Fehler werden erkannt: falsche Belegungsnorm, verwechselte lokale Architektur, verkürzter Energiefilter, verlorene Resetwahrscheinlichkeit, defekte Ginsparg--Wilson-Relation, entfernte CAR-Kreuzungszeichen und falsches Spektralpolynom. Die meisten Bedingungen sind einzelne Zustände oder Übergänge, nicht unabhängige physikalische Entdeckungen.

Separat wurde der 24.024-dimensionale Singulettoperator neu aufgebaut und sein korrigierter niedriger Spektralbereich berechnet. Dieser Lauf ist numerisch, keine formale Eigenwertzertifizierung und kein erneut im Optimierungsmodus ausgeführter Großlauf. Die exakten Argumente sind ausgeschrieben und durch endliche Kontrollen ergänzt, nicht als Lean-Beweise formalisiert. Forschungsbericht, Rohwerte, Prüfer, unveränderte Eingaben und Dokumentquellen liegen im zugehörigen versionierten Paket. Ältere Befunde behalten ihre ursprüngliche Provenienz; die historische Prüfanzahl wird nicht als neuer Lauf ausgegeben.

\begin{grenze}{Abschluss dieser Revision, nicht Abschluss der TOE}
Der stärkste neue positive Satz betrifft ein vollständig bezahltes endliches kontrolliertes Experiment. Die stärkste Vereinfachung betrifft die positive lokale F4-Darstellung. Die wichtigste offene Frage ist die gemeinsame Herkunft dieser Struktur und ihrer Kontrollen aus dem Compiler. Kein vollständiges T1--T8-Tor, kein RH-Beweis, kein neues allgemeines Faktorisierungsverfahren und keine P-versus-NP-Lösung wird als geschlossen markiert.
\end{grenze}
